% ============================================================
% Durūs al-Muhimmah – Unterrichtsausgabe
% Gemeinsame Präambel: Sprachen, Schriften, Farben, Umgebungen
% ============================================================
% WICHTIG: Das bidi-Paket (via polyglossia + \setotherlanguage{arabic})
% verlangt, ganz am Ende geladen zu werden -- nach amsmath, xcolor,
% graphicx, tikz, tcolorbox, fancyhdr, hyperref usw. Reihenfolge nicht
% ohne Grund ändern.

\usepackage{fontspec}
\usepackage{amsmath}
\usepackage[a4paper,inner=14mm,outer=10mm,top=12mm,bottom=12mm,bindingoffset=4mm]{geometry}
\usepackage{xcolor}
\usepackage{graphicx}
\usepackage{tikz}
\usetikzlibrary{calc}

% Farbpalette – extrahiert aus dem DEEN-Bildungszentrum-Logo und dem
% persönlichen Kunya-Logo (Abu 'A'isha al-Kumaisi) von Sahdjad Mahamadou.
% Nicht per Hand ändern, ohne EDITORIAL_NOTES.md zu aktualisieren.

\definecolor{PrimaryGreen}{HTML}{4F8A3D}      % Hauptgrün (Logo-Mittelton, Buchdeckel/Ikonen)
\definecolor{DarkGreen}{HTML}{123018}          % dunkles Waldgrün (Überschriften, Kunya-Logo-Hintergrund)
\definecolor{DeepestGreen}{HTML}{081209}       % fast Schwarzgrün (Cover-Hintergrund, Lektionstrennseiten)
\definecolor{SecondaryGreen}{HTML}{6FA85A}     % Salbeigrün (Akzente, Tabellenköpfe)
\definecolor{LightSage}{HTML}{9FC98A}          % helles Blattgrün (dezente Flächen, Muster)
\definecolor{LightBackground}{HTML}{F3F7EF}    % helle Hintergrundfläche für Boxen
\definecolor{OffWhite}{HTML}{FBFAF6}           % Seitenhintergrund / Buchblock
\definecolor{NeutralText}{HTML}{1E1E1E}        % Fließtext
\definecolor{MutedText}{HTML}{5A5A52}          % Kopf-/Fußzeilen, sekundärer Text
\definecolor{Accent}{HTML}{B08D3E}             % warmer Goldton für seltene Hervorhebung (Merke-Rahmen)

% Lektionsakzente (Abstufungen von Grün, siehe Master-Spezifikation §22)
\definecolor{Lesson01}{HTML}{123018}
\definecolor{Lesson02}{HTML}{2C4A22}
\definecolor{Lesson03}{HTML}{3D5C2A}
\definecolor{Lesson04}{HTML}{4F8A3D}
\definecolor{Lesson05}{HTML}{1F5C3F}
\definecolor{Lesson06}{HTML}{2E6B4A}
\definecolor{Lesson07}{HTML}{35785A}
\definecolor{Lesson08}{HTML}{4F8A3D}
\definecolor{Lesson09}{HTML}{6FA85A}
\definecolor{Lesson10}{HTML}{2C4A22}
\definecolor{Lesson11}{HTML}{123018}
\definecolor{Lesson12}{HTML}{3D5C2A}
\definecolor{Lesson13}{HTML}{1F5C3F}
\definecolor{Lesson14}{HTML}{4F8A3D}
\definecolor{Lesson15}{HTML}{2E6B4A}
\definecolor{Lesson16}{HTML}{35785A}
\definecolor{Lesson17}{HTML}{6FA85A}
\definecolor{Lesson18}{HTML}{123018}

\graphicspath{{\PROJROOT/assets/}}
\usepackage[most]{tcolorbox}
\usepackage{fancyhdr}
\usepackage{enumitem}
\usepackage{tabularx}
\usepackage{booktabs}
\usepackage{needspace}
\usepackage{titlesec}
\usepackage{eso-pic}
\usepackage[colorlinks=true,linkcolor=DarkGreen,urlcolor=DarkGreen,citecolor=DarkGreen,bookmarksopen=true]{hyperref}
\usepackage{bookmark}
\usepackage[none]{hyphenat}
\usepackage{ragged2e}

% -- Sprachen (bidi/arabic ganz zuletzt laden) -------------------
\usepackage{polyglossia}
\setdefaultlanguage{german}
\setotherlanguage{arabic}

% -- Schriften -------------------------------------------------
% Bewusst schlichte, "computerhafte" Druckschriften statt kalligrafisch
% wirkender Fonts (Nutzerfeedback: vorherige Wahl sah zu sehr nach
% Handschrift aus).
\setmainfont{Latin Modern Roman}
\newfontfamily\headingfont{TeX Gyre Heros}[Scale=0.95]
\newfontfamily\arabicfont[Script=Arabic]{Noto Naskh Arabic}
\newfontfamily\quranfont[Script=Arabic]{Noto Naskh Arabic}

\newcommand{\ar}[1]{{\arabicfont\begin{Arabic}#1\end{Arabic}}}
\newcommand{\quran}[1]{{\quranfont\begin{Arabic}#1\end{Arabic}}}

\setlength{\parindent}{0pt}
\setlength{\parskip}{0.6\baselineskip}
\linespread{1.12}

% -- Silbentrennung am Zeilenende komplett deaktivieren (Nutzerwunsch:
%    keine Bindestriche im Buch, außer dort, wo sie inhaltlich hingehören,
%    z. B. in Transliterationen wie "al-Fātiḥa") ----------------------
\hyphenpenalty=10000
\exhyphenpenalty=10000
\RaggedRight

% -- Überschriften: GROSSBUCHSTABEN statt Groß-/Kleinschreibung -------
\titleformat{\section}{\headingfont\LARGE\bfseries\color{DarkGreen}}{}{0pt}{\MakeUppercase}
\titleformat{\subsection}{\headingfont\Large\bfseries\color{PrimaryGreen}}{}{0pt}{\MakeUppercase}
\titleformat{\subsubsection}{\headingfont\large\bfseries\color{SecondaryGreen}}{}{0pt}{\MakeUppercase}
\titlespacing*{\section}{0pt}{1.4em}{1em}
\titlespacing*{\subsection}{0pt}{1.2em}{0.8em}

% -- Kopf-/Fußzeile ----------------------------------------------
\pagestyle{fancy}
\fancyhf{}
\renewcommand{\headrulewidth}{0.4pt}
\renewcommand{\footrulewidth}{0pt}
\fancyhead[LE]{\small\color{MutedText}\textsc{Durūs al-Muhimmah}}
\fancyhead[RO]{\small\color{MutedText}\nouppercase{\leftmark}}
\fancyfoot[LE,RO]{\small\color{MutedText}\thepage}

\fancypagestyle{plain}{%
  \fancyhf{}
  \fancyfoot[LE,RO]{\small\color{MutedText}\thepage}
  \renewcommand{\headrulewidth}{0pt}
}

\fancypagestyle{leer}{%
  \fancyhf{}
  \renewcommand{\headrulewidth}{0pt}
}

% ============================================================
% Didaktische Umgebungen
% ============================================================

% -- "Der Shaykh رحمه الله sagt:" -------------------------------
\newtcolorbox{shaykhsagt}{
  enhanced, breakable,
  colback=LightBackground, colframe=DarkGreen,
  boxrule=0.9pt, arc=2pt,
  left=12pt, right=12pt, top=10pt, bottom=10pt,
  before skip=14pt, after skip=14pt,
  fonttitle=\headingfont\bfseries\color{OffWhite},
  colbacktitle=DarkGreen,
  title={Der Shaykh \ar{رحمه الله} sagt:}
}

% -- Lernziele ----------------------------------------------------
\newtcolorbox{lernziele}{
  enhanced, breakable,
  colback=OffWhite, colframe=SecondaryGreen,
  boxrule=0.8pt, arc=2pt,
  left=12pt, right=12pt, top=8pt, bottom=8pt,
  before skip=14pt, after skip=14pt,
  fonttitle=\headingfont\bfseries\color{DarkGreen},
  title={Nach dieser Lektion solltest du \dots}
}

% -- Merke-Kasten ---------------------------------------------------
\newtcolorbox{merke}{
  enhanced, breakable,
  colback=LightSage!35!OffWhite, colframe=Accent,
  boxrule=1pt, arc=2pt,
  left=12pt, right=12pt, top=8pt, bottom=8pt,
  before skip=12pt, after skip=12pt,
  fonttitle=\headingfont\bfseries\color{DarkGreen},
  title={MERKE}
}

% -- Beispiel ------------------------------------------------------
\newtcolorbox{beispiel}{
  enhanced, breakable,
  colback=OffWhite, colframe=PrimaryGreen,
  boxrule=0.6pt, arc=2pt, dashed,
  left=12pt, right=12pt, top=8pt, bottom=8pt,
  before skip=12pt, after skip=12pt,
  fonttitle=\headingfont\bfseries\color{PrimaryGreen},
  title={Beispiel}
}

% -- Hinweis ---------------------------------------------------------
\newtcolorbox{hinweis}{
  enhanced, breakable,
  colback=OffWhite, colframe=MutedText,
  boxrule=0.6pt, arc=2pt,
  left=12pt, right=12pt, top=8pt, bottom=8pt,
  before skip=12pt, after skip=12pt,
  fonttitle=\headingfont\bfseries\color{MutedText},
  title={Hinweis}
}

% -- Zusammenfassung --------------------------------------------------
\newtcolorbox{zusammenfassung}{
  enhanced, breakable,
  colback=DarkGreen!6!OffWhite, colframe=DarkGreen,
  boxrule=0.9pt, arc=2pt,
  left=12pt, right=12pt, top=10pt, bottom=10pt,
  before skip=16pt, after skip=10pt,
  fonttitle=\headingfont\bfseries\color{OffWhite},
  colbacktitle=DarkGreen,
  title={Das Wichtigste zusammengefasst}
}

% -- Qur'an-Vers-Block --------------------------------------------------
\newenvironment{quranvers}[1]{%
  \def\quranversref{#1}%
  \begin{tcolorbox}[enhanced,breakable,colback=LightBackground,colframe=SecondaryGreen,
    boxrule=0.7pt,arc=2pt,left=12pt,right=12pt,top=8pt,bottom=8pt,
    before skip=10pt,after skip=10pt]
  \begin{center}
}{%
  \end{center}
  \hfill{\small\color{MutedText}[\quranversref]}
  \end{tcolorbox}
}

% -- Begriffe-Tabelle ------------------------------------------------
\newenvironment{begriffe}{%
  \subsubsection*{Wichtige Begriffe dieser Lektion}
  \renewcommand{\arraystretch}{1.4}
  \begin{tabularx}{\linewidth}{@{}l >{\raggedright\arraybackslash}X@{}}
  \toprule
  \textbf{\color{DarkGreen}Begriff} & \textbf{\color{DarkGreen}Kurze Bedeutung} \\
  \midrule
}{%
  \bottomrule
  \end{tabularx}
}

% -- Wiederkehrendes Ornament-Muster (eigene Gestaltung, farblich an das
%    Kunya-Logo angelehnt, keine 1:1-Kopie) --------------------------------
% #1 = Basisfarbe (Vollflächenfüllung), #2 = Mischanteil der Musterfarbe in
% Prozent (0..100), als deckende Mischfarbe berechnet -- bewusst KEIN
% TikZ/pgf "opacity=", weil das zusammen mit eso-pic-Shipout-Hintergründen
% eine defekte ExtGState-Ressource erzeugt (PDF wird unlesbar / Inhalt nach
% dem Hintergrund verschwindet). Nicht wieder auf opacity umstellen.
\newcommand{\ornamentmuster}[2]{%
  \begin{tikzpicture}
    \fill[#1] (0mm,0mm) rectangle (\paperwidth,\paperheight);
    \begin{scope}
      \clip (0mm,0mm) rectangle (\paperwidth,\paperheight);
      \foreach \row in {0,...,9}{%
        \foreach \col in {0,...,6}{%
          \pgfmathsetmacro{\xoff}{ (mod(\row,2))*17 }
          \coordinate (P) at (\col*34mm+\xoff mm+10mm,\row*33mm+8mm);
          \draw[CoverPatternA!#2!#1,line width=0.55pt]
            (P) ++(0mm,-9mm) .. controls +(6mm,-2mm) and +(6mm,2mm) .. ++(0mm,9mm)
                .. controls +(-6mm,2mm) and +(-6mm,-2mm) .. ++(0mm,-9mm);
          \draw[CoverPatternB!#2!#1,line width=0.55pt]
            (P) ++(-6mm,0mm) circle[radius=1.1mm];
          \draw[AccentLight!#2!#1,line width=0.45pt]
            (P) ++(0mm,9mm) -- ++(0mm,4mm);
          \draw[CoverPatternA!#2!#1,line width=0.45pt]
            (P) ++(9mm,-4.5mm) arc[start angle=0,end angle=180,radius=9mm];
        }
      }
    \end{scope}
  \end{tikzpicture}%
}

% -- Lektionstrennseite -------------------------------------------------
\newcommand{\lektionstrennseite}[3]{%
  % #1 = Nummer (arabisch, z.B. 01), #2 = römisch/ausgeschrieben für Farbwahl, #3 = Titel
  \clearpage
  \thispagestyle{empty}
  \AddToShipoutPictureBG*{\ornamentmuster{Lesson#2}{55}}
  \null\vspace*{0.30\textheight}
  \begin{center}
    {\headingfont\color{AccentLight}\fontsize{22}{26}\selectfont LEKTION}\\[2mm]
    {\headingfont\bfseries\color{OffWhite}\fontsize{72}{80}\selectfont #1}\\[10mm]
    {\headingfont\bfseries\color{OffWhite}\fontsize{20}{24}\selectfont \MakeUppercase{#3}}
  \end{center}
  \clearpage
}
